\chapter{Perímetro, área y volumen}\label{ch:g6:measure}

¿Cuánto mide la cerca, qué tamaño tiene el campo, cuánta agua cabe en el
tanque? Tres preguntas diferentes, tres cantidades diferentes ---
\omterm{def:g6:measure:perimeter}{perímetro}, \omterm{def:g6:measure:area}{área}, \omterm{def:g6:measure:volume}{volumen} --- cada uno con sus propias unidades. Confundirlos es
el error más común en geometría; este capítulo los clasifica para
bueno.

\section{Longitudes y perímetro}

\begin{definition}[Perímetro]\label{def:g6:measure:perimeter}
El \emph{perímetro}\index{perímetro} de una figura es la longitud total
de su frontera. Se mide en unidades de longitud: milímetros (mm),
centímetros (cm), metros (m), kilómetros (km), con
\[
1 \text{ km} = 1000 \text{ m}, \qquad
1 \text{ m} = 100 \text{ cm}, \qquad
1 \text{ cm} = 10 \text{ mm}.
\]
\end{definition}

\begin{example}\label{ex:g6:measure:perimeter}
Un rectángulo de longitud $L $ y ancho $ w $ tiene \omterm{def:g6:measure:perimeter}{perímetro}
\[
P = L + w + L + w = 2 \times (L + w).
\]
por un $7$ centímetros por $4$ cm rectángulo: $ P = 2 \times (7 + 4) = 2 \times 11 =
22$ centímetro. un cuadrado de lado $ c $ tiene \omterm{def:g6:measure:perimeter}{perímetro} $4c $.
\end{example}

\begin{proposition}[Circunferencia de un círculo]\label{prop:g6:measure:circle}
El \omterm{def:g6:measure:perimeter}{perímetro} (o \emph{circunferencia}) de un círculo de \omterm{def:g3:shapes:shapes}{radio} $ r $ es
\[
P = 2 \pi r ,
\]
dónde $\pi \approx 3.14$ es el mismo número para cada círculo.
\end{proposition}
\admitted

\begin{example}
Un estanque circular tiene \omterm{def:g3:shapes:shapes}{radio} $5$ metro. Sus medidas fronterizas
$2 \times \pi \times 5 = 10\pi \approx 31.4$ metro. Mantener el valor exacto
$10\pi $ el mayor tiempo posible; \omterm{def:g4:numbers:round}{redondo} sólo al final.
\end{example}

\section{Áreas}

\begin{definition}[Área]\label{def:g6:measure:area}
El \emph{área}\index{área} de una figura mide la superficie que cubre:
¿Cuántos cuadrados unitarios caben dentro? Unidades: cm$^2$ (un cuadrado de lado
$1$ centímetros), metros $^2$, kilómetros $^2$, \dots\ Cuidadoso:
\[
1 \text{ m}^2 = 100 \times 100 \text{ cm}^2 = 10\,000 \text{ cm}^2
\]
--- un metro cuadrado es un $100$ centímetros por $100$ cm cuadrados, por lo que cada paso de
la escalera de las unidades vale $100$, no $10$.
\end{definition}

\begin{omfigure}
\begin{tikzpicture}[scale=0.68]
  \draw[gray!50, thin] (0,0) grid (7,4);
  \draw[very thick, omDef] (0,0) rectangle (7,4);
  \fill[omThm!25] (0,3) rectangle (1,4);
  \draw[omThm, thick] (0,3) rectangle (1,4);
  \node[font=\small, omThm] at (2.6,3.5) {$1$ unit square};
  \node[font=\small] at (3.5,-0.5) {$7 \times 4 = 28$ unit squares};
\end{tikzpicture}

{\small El \omterm{def:g6:measure:area}{área} de un rectángulo: $7$ columnas de $4$ cuadrados unitarios cada uno,
entonces $7 \times 4 = 28$ cuadrados en total. Esta es la razón \omterm{def:g6:measure:area}{área} $=$ longitud
$\times $ ancho.}
\end{omfigure}

\begin{proposition}[Fórmulas de área básicas]\label{prop:g6:measure:formulas}
\begin{center}
\renewcommand{\arraystretch}{1.5}
\begin{tabular}{l|l}
figure & \omterm{def:g6:measure:area}{área} \\
\hline
rectangle ($ L \times w $) & $ A = L \times w $ \\
square (side $ c $) & $ A = c^2$ \\[4pt]
\omterm{def:g6:shapes:triangles}{triángulo rectángulo} (legs $ a $, $ b $) & $ A = \dfrac{a \times b}{2}$ \\[4pt]
disk (\omterm{def:g3:shapes:shapes}{radio} $ r $) & $ A = \pi r^2$ \\
\end{tabular}
\end{center}
\end{proposition}

\begin{proof}[Prueba del triángulo rectángulo]
Dos copias de un \omterm{def:g6:shapes:triangles}{triángulo rectángulo} con piernas $ a $ y $ b $, pegado a lo largo del
hipotenusa, forma una $ a \times b $ rectángulo. Entonces el triángulo \omterm{def:g6:measure:area}{área} es
\omterm{def:g3:division:half}{medio} el rectángulo: $\frac{ab}{2}$. (La fórmula del rectángulo es la
conteo de unidades cuadradas arriba; se admite la fórmula del disco, ver
\cref{ch:g7:areas}.)
\end{proof}

\begin{omfigure}
\begin{tikzpicture}[scale=0.75]
  \draw[thick, fill=omDef!15] (0,0) -- (4,0) -- (0,2.5) -- cycle;
  \draw[thick, omExo, dashed] (4,0) -- (4,2.5) -- (0,2.5);
  \draw[thin] (0.3,0) -- (0.3,0.3) -- (0,0.3);
  \node[below, font=\small] at (2,-0.1) {$ a $};
  \node[left, font=\small] at (0,1.25) {$ b $};
  \node[font=\small, omExo] at (2.8,1.7) {second copy};
\end{tikzpicture}

{\small A \omterm{def:g6:shapes:triangles}{triángulo rectángulo} es \omterm{def:g3:division:half}{medio} un rectángulo: de ahí la fórmula
$\frac{a\times b}{2}$.}
\end{omfigure}

\begin{example}[Figuras compuestas]\label{ex:g6:measure:composite}
Una habitación en forma de L es una $6$ metro $\times $ $4$ m rectángulo con un $2$ metro
$\times $ $2$ m esquina cuadrada eliminada. Es \omterm{def:g6:measure:area}{área}, paso a paso:
\begin{enumerate}
  \item rectángulo completo: $6 \times 4 = 24$ metro $^2$;
  \item cuadrado eliminado: $2 \times 2 = 4$ metro $^2$;
  \item restante \omterm{def:g6:measure:area}{área}: $24 - 4 = 20$ metro $^2$.
\end{enumerate}
Es \emph{\omterm{def:g6:measure:perimeter}{perímetro}} no es $20$ cualquier cosa: caminar por la L, la
La frontera aún mide $6 + 4 + 6 + 4 = 20$ m --- los dos cortes de la
esquina reemplace dos piezas iguales de pared. Mismo número por coincidencia,
¡Pero metros cuadrados para uno, metros para el otro!
\end{example}

\begin{remark}[Mismo perímetro, diferentes zonas]
Dos figuras pueden tener el mismo \omterm{def:g6:measure:perimeter}{perímetro} y muy diferente \omterm{def:g6:measure:area}{áreas}: a
$5 \times 5$ cuadrado y un $9 \times 1$ rectángulo ambos tienen \omterm{def:g6:measure:perimeter}{perímetro}
$20$, pero \omterm{def:g6:measure:area}{áreas} $25$ y $9$. \omterm{def:g6:measure:perimeter}{Perímetro} y \omterm{def:g6:measure:area}{área} son verdaderamente independientes
cantidades.
\end{remark}

\section{Volúmenes}

\begin{definition}[Volumen]\label{def:g6:measure:volume}
El \emph{volumen}\index{volumen} de un \omterm{def:g5:solids:def}{sólido} mide el espacio que ocupa:
cuantas unidades \omterm{def:g5:solids:def}{cubos} encajar en el interior. Unidades: cm $^3$, metro $^3$, \dots\ con
$1$ metro $^3 = 1\,000\,000$ centímetro $^3$ (cada paso de la escalera vale
$1000$). Para líquidos también se utiliza el \emph{litro}:
\[
1 \text{ L} = 1 \text{ dm}^3 = 1000 \text{ cm}^3,
\qquad
1 \text{ m}^3 = 1000 \text{ L}.
\]
\end{definition}

\begin{proposition}[volumen de una caja]\label{prop:g6:measure:box}
un rectangular \omterm{def:g5:solids:def}{caja} (a \emph{prisma rectangular}) de longitud $ L $, ancho $ w $
y altura $ h $ tiene \omterm{def:g6:measure:volume}{volumen}
\[
V = L \times w \times h .
\]
\end{proposition}

\begin{proof}
La capa inferior contiene $ L \times w $ unidad \omterm{def:g5:solids:def}{cubos}
(\cref{def:g6:measure:area} imagen, con \omterm{def:g5:solids:def}{cubos}), y hay $ h $
tales capas.
\end{proof}

\begin{omfigure}
\begin{tikzpicture}[scale=0.6]
  \foreach \y in {0,1} {
    \foreach \x in {0,1,2,3} {
      \foreach \z in {0,1} {
        \begin{scope}[shift={({\x+0.45*\z},{\y+0.3*\z})}]
          \draw[fill=omDef!12] (0,0) rectangle (1,1);
          \draw[fill=omDef!20] (0,1) -- (0.45,1.3) -- (1.45,1.3) -- (1,1)
            -- cycle;
          \draw[fill=omDef!28] (1,0) -- (1.45,0.3) -- (1.45,1.3) -- (1,1)
            -- cycle;
        \end{scope}
      }
    }
  }
  \node[below, font=\small] at (2.5,-0.3)
    {$4 \times 2 \times 2 = 16$ unit cubes};
\end{tikzpicture}

{\small A \omterm{def:g5:solids:def}{caja} lleno de unidad \omterm{def:g5:solids:def}{cubos}: $4 \times 2$ \omterm{def:g5:solids:def}{cubos} por capa, dos
capas.}
\end{omfigure}

\begin{example}\label{ex:g6:measure:aquarium}
Un acuario mide $60$ centímetros por $30$ centímetros por $40$ cm (altura). \omterm{def:g6:measure:volume}{Volumen}:
\[
V = 60 \times 30 \times 40 = 72\,000 \text{ cm}^3 = 72 \text{ dm}^3
= 72 \text{ L}.
\]
Paso a paso para la conversión: $1000$ centímetro $^3$ hacer un litro, y
$72\,000 \div 1000 = 72$.
\end{example}

\section{Ejercicios}

\begin{exercise}[$\star $]\label{exo:g6:measure:1}
Convertir: $3.5$ m en cm; $420$ mm en cm; $0.8$ km en
metro;
$25\,000$ m en km.
\end{exercise}

\begin{exercise}[$\star $]\label{exo:g6:measure:2}
Calcular el \omterm{def:g6:measure:perimeter}{perímetro} de: un rectángulo $8$ centímetro $\times $ $3.5$ centímetro; un cuadrado
de lado $6.2$ centímetro; un triángulo con lados $5$ centímetro, $7$ centímetros y
$9$ centímetro.
\end{exercise}

\begin{exercise}[$\star $]\label{exo:g6:measure:3}
Una pista circular para correr tiene \omterm{def:g3:shapes:shapes}{radio} $50$ metro. cuanto dura una vuelta
(valor exacto con $\pi $, luego redondeado al metro)? cuantas vueltas dan
al menos $2$ kilómetros?
\end{exercise}

\begin{exercise}[$\star $]\label{exo:g6:measure:4}
Calcular el \omterm{def:g6:measure:area}{área} de: un rectángulo $9$ centímetro $\times $ $4$ centímetro; un cuadrado
de lado $7$ metro; a \omterm{def:g6:shapes:triangles}{triángulo rectángulo} con piernas $6$ centímetros y $10$ centímetro; un disco de
\omterm{def:g3:shapes:shapes}{radio} $3$ cm (valor exacto, luego redondeado a cm $^2$).
\end{exercise}

\begin{exercise}[$\star $]\label{exo:g6:measure:5}
Convertir: $3$ metro $^2$ en centímetros $^2$; $45\,000$ centímetro $^2$ en m $^2$; $2.5$ l
en centímetros $^3$; $4\,500$ l en m $^3$.
\end{exercise}

\begin{exercise}[$\star $]\label{exo:g6:measure:6}
Un campo rectangular es $120$ m de largo y $85$ m de ancho. cuantos metros de
¿Se necesita una valla para encerrarlo? cual es su \omterm{def:g6:measure:area}{área}?
\end{exercise}

\begin{exercise}[$\star $]\label{exo:g6:measure:7}
Calcular el \omterm{def:g6:measure:volume}{volumen} de un \omterm{def:g5:solids:def}{caja} $5$ centímetro $\times $ $4$ centímetro $\times $ $10$ centímetros,
y de un \omterm{def:g5:solids:def}{cubo} de borde $3$ centímetro.
\end{exercise}

\begin{exercise}[$\star $]\label{exo:g6:measure:8}
Dibuja dos rectángulos diferentes con \omterm{def:g6:measure:perimeter}{perímetro} $16$ cm y calcular su
\omterm{def:g6:measure:area}{áreas}. ¿Cuál de tus rectángulos tiene el mayor \omterm{def:g6:measure:area}{área}?
\end{exercise}

\begin{exercise}[$\star\star $]\label{exo:g6:measure:9}
Una figura en forma de T está hecha de un $10 \times 2$ rectángulo horizontal en
parte superior de un $2 \times 6$ vertical (medidas en cm). Calcular su
\omterm{def:g6:measure:area}{área}, entonces es \omterm{def:g6:measure:perimeter}{perímetro} (camina alrededor de la frontera con cuidado, agregando
cada borde).
\end{exercise}

\begin{exercise}[$\star\star $]\label{exo:g6:measure:10}
Una piscina es una forma rectangular. \omterm{def:g5:solids:def}{caja} $25$ m de largo, $10$ m de ancho y
$2$ m de profundidad.
\begin{enumerate}
  \item ¿Cuántos metros cúbicos de agua cabe cuando está lleno?
  \item ¿Cuantos litros es eso?
  \item La piscina se llena en $50\,000$ L por hora. ¿Cuánto dura el
        llenado tomar?
\end{enumerate}
\end{exercise}

\begin{exercise}[$\star\star $]\label{exo:g6:measure:11}
Un jardín es un cuadrado de lado. $20$ m que contiene un estanque circular de
\omterm{def:g3:shapes:shapes}{radio} $3$ metro. Qué \omterm{def:g6:measure:area}{área} de pasto hay (valor exacto, luego redondeado)
a la m $^2$)?
\end{exercise}

\begin{exercise}[$\star\star\star $]\label{exo:g6:measure:12}
Una barra de chocolate mide $15$ centímetro $\times $ $6$ centímetro $\times $ $1$ centímetro.
El fabricante se duplica \emph{los tres} dimensiones.
\begin{enumerate}
  \item ¿A cuánto asciende el \omterm{def:g6:measure:volume}{volumen} ¿multiplicado? Verificar calculando ambos
        \omterm{def:g6:measure:volume}{volúmenes}.
  \item El precio se multiplica por $4$. ¿Es el gran bar una mejor oferta?
        que el pequeño?
\end{enumerate}
\end{exercise}

\section{Problema: la valla de la reina Dido}

\begin{problem}[{Problema de fin de semana --- Con una longitud fija de cerca, ¿qué forma encierra la mayor cantidad de terreno? El cuadrado vence a cada rectángulo, el círculo vence al cuadrado y el muro de un granero lo cambia todo}]
\label{pb:g6:measure:1}
La leyenda dice que la reina Dido, al desembarcar en la costa de África, fue
se le concedió "tanta tierra como la que puede abarcar una piel de buey", por lo que cortó
la piel en una tira inmensamente larga y delgada y lo suficientemente encerrada
terreno para fundar la ciudad de Cartago. Su problema ahora es tuyo: una
el granjero posee exactamente $20$ m de valla. \Cref{exo:g6:measure:8}
demostró que dos plumas con el mismo \omterm{def:g6:measure:perimeter}{perímetro} puede tener diferentes
\omterm{def:g6:measure:area}{áreas}; este problema encuentra el \emph{mejor} pluma --- y descubre
que la respuesta cambia por completo cuando la pared de un granero da acceso gratuito
lado.

\medskip
\noindent\textbf{Parte I --- Twenty meters of fence.}
\begin{enumerate}
  \item El bolígrafo debe ser un rectángulo usando todos $20$ m de valla.
        Compruebe que un $1 \times 9$ bolígrafo y un $2 \times 8$ pluma ambos
        calificar y calcular su \omterm{def:g6:measure:area}{áreas}.
  \item Explique por qué el largo y el ancho de cada corral calificado
        sumar hasta $10$ metro. Luego haz la tabla completa de
        bolígrafos de números enteros ($1 \times 9$ arriba a $5 \times 5$) con
        su \omterm{def:g6:measure:area}{áreas}. ¿Cuál es mejor?
  \item ¿Vale la pena probar los lados decimales? Calcular el \omterm{def:g6:measure:area}{áreas} del
        $4.5 \times 5.5$ pluma y de la $4.9 \times 5.1$ bolígrafo, y
        comparar con el $5 \times 5$ cuadrado. ¿Qué haces?
        conjetura?
  \item La pregunta 2 convirtió el problema de la valla en un número puro
        pregunta: \emph{entre todos los pares de números que suman
        $10$, ¿qué par tiene el mayor \omterm{def:g2:mult:def}{producto}?} Respondelo desde
        tu tabla y establece la regla general que observas.
  \item He aquí por qué alejarse de la plaza siempre pierde,
        con tijeras en lugar de álgebra: empezar desde el
        $5 \times 5$ bolígrafo y cámbielo al $6 \times 4$ bolígrafo
        quitando una tira y pegando otra hacia atrás. cual
        se quita la tira, cuál se agrega y por qué se quita la tira
        ¿El intercambio pierde exactamente un metro cuadrado? Explique por qué un
        paso adicional (para $7 \times 3$) pierde \omterm{def:g3:numbers:evenodd}{incluso} más.
\end{enumerate}

\medskip
\noindent\textbf{Parte II --- The barn, and the circle.}
\begin{enumerate}[resume]
  \item El corral ahora está construido contra una larga pared del granero: la pared
        reemplaza una longitud del bolígrafo, y el $20$ m de valla
        cubrir sólo los otros tres lados. Haz una tabla de
        bolígrafos de números enteros (ancho $ w $, longitud $ L $ a lo largo de la pared,
        $ w + L + w = 20$) y sus \omterm{def:g6:measure:area}{áreas}. ¿Qué bolígrafo gana ahora? ---
        y es un cuadrado?
  \item Explica el ganador con un espejo.
        (\cref{ch:g6:symmetry}): refleja el bolígrafo a través del granero
        pared y considere el corral duplicado. ¿Cuánto cuesta la cerca?
        el uso de pluma duplicada, qué pluma duplicada es mejor según la Parte I,
        ¿Y qué hace eso con el bolígrafo original?
  \item Dido no construyó un rectángulo. Doblar el $20$ m de valla
        en un círculo: usando $\pi \approx 3.14$, calcular su
        \omterm{def:g3:shapes:shapes}{radio} (al cm), entonces es \omterm{def:g6:measure:area}{área} (a la m $^2$)
        (\cref{prop:g6:measure:circle},
        \cref{prop:g6:measure:formulas}). Comparar con el
        $5 \times 5$ cuadrado.
  \item El problema inverso: un bolígrafo de \omterm{def:g6:measure:area}{área} exactamente $36$ metro $^2$ es
        deseado, con la menor valla posible. comparar el
        rectángulos $1 \times 36$, $2 \times 18$, $3 \times 12$,
        $4 \times 9$ y $6 \times 6$: \omterm{def:g6:measure:perimeter}{perímetros}? ¿Qué forma es
        más barato, y ¿cómo es que esta pregunta es el reflejo de
        ¿Parte I?
  \item En una frase: ¿por qué las latas, las tuberías y los tanques de agua son tan
        a menudo \emph{\omterm{def:g4:numbers:round}{redondo}}?
\end{enumerate}

\medskip
\noindent\textbf{Parte III --- \omterm{def:g6:measure:perimeter}{Perímetro} and \omterm{def:g6:measure:area}{área} are strangers.}
\begin{enumerate}[resume]
  \item Encuentra un rectángulo cuyo \omterm{def:g6:measure:perimeter}{perímetro} excede $100$ m pero
        cuyo \omterm{def:g6:measure:area}{área} es menor que $1$ metro $^2$. (Muy largo y muy
        delgado. Se permiten decimales.)
  \item Verdadero o falso: ``de dos figuras, la de mayor
        \omterm{def:g6:measure:perimeter}{perímetro} tiene el mas grande \omterm{def:g6:measure:area}{área}." Da un contraejemplo
        de este problema.
  \item Toma el $6 \times 4$ rectángulo y cortar un $1 \times 1$
        muesca cuadrada en el medio de un lado largo (la muesca
        se abre hacia afuera, como un diente perdido). Calcular el \omterm{def:g6:measure:area}{área}
        y el \omterm{def:g6:measure:perimeter}{perímetro} de la figura con muescas y compare ambos
        con el original. ¿Qué pasó con cada uno?
  \item Los cartógrafos conocen un hecho extraño: cuanto más detallado es el mapa,
        el \emph{más extenso} una costa mide --- cada zoom
        revela nuevas pequeñas muescas, y la pregunta 13 muestra lo que
        cada muesca lo hace. Explica en una o dos frases por qué.
        ``la longitud de la costa de Bretaña'' es una superficie resbaladiza
        número, mientras que "el \omterm{def:g6:measure:area}{área} de Bretaña'' no lo es.
  \item El examen del granjero. Con $36$ m de valla y el granero
        pared disponible, encuentre el mejor bolígrafo rectangular (use el
        espejo de la pregunta 7), dé su \omterm{def:g6:measure:area}{área}y comparar con
        el mejor bolígrafo construido con $36$ m en campo abierto. cuanto
        ¿La pared del granero le gana al granjero?
\end{enumerate}
\end{problem}
